\subsection{EXISTENCE AND UNIQUENESS OF SOLUTIONS OF LIPSCHITZ AND LOCALLY LIPSCHITZ EQUATIONS}

THEOREM 1 (Cauchy). Let \(\mathbf{f}\) be a Lipschitz function on \(\mathrm{I} \times \mathrm{H}\) , let \(\mathrm{J}\) be a compact subinterval of \(\mathrm{I}\) , not reducing to a single point, \(t_0\) a point of \(\mathrm{J}\) , \(\mathrm{S}\) an open ball with centre \(\mathbf{x}_0\) and radius \(r\) , contained in \(\mathrm{H}\) , and \(\mathrm{M}\) the least upper bound of \(\| \mathbf{f}(t, \mathbf{x})\|\) on \(\mathrm{J} \times \mathrm{S}\) . In these circumstances, for every compact interval \(\mathrm{K}\) that does not reduce to a single point and is contained in the intersection of \(\mathrm{J}\) with \(]t_0 - r / M\) , \(t_0 + r / M[\) , and contains \(t_0\) , there exists one and only one solution of the differential equation \(\mathbf{x}' = \mathbf{f}(t, \mathbf{x})\) defined on \(\mathrm{K}\) , with values in \(\mathrm{S}\) , and equal to \(\mathbf{x}_0\) at the point \(t_0\) .

Indeed, for \(\varepsilon > 0\) sufficiently small, the set \(F_{\varepsilon}\) of approximate solutions of (1) to within \(\varepsilon\) , defined on K, with values in S, and equal to \(x_{0}\) at the point \(t_{0}\) , is not empty (IV, p.~166, prop. 3); further, if u and v belong to \(F_{\varepsilon}\) then, by (15) (IV, p.~170)

\[
\| \mathbf {u} (t) - \mathbf {v} (t) \| \leqslant 2 \varepsilon \frac {e ^ {k | t - t _ {0} |} - 1}{k}
\]

for all \(t \in K\) , so the sets \(F_{\varepsilon}\) form a filter base G which converges uniformly on K to a continuous function w, equal to \(x_{0}\) at \(t_{0}\) ; also w takes values in S, since, for \(\varepsilon\) small enough, the functions \(u \in F_{\varepsilon}\) take their values in a closed ball contained in S. Since \(\mathbf{f}(t, \mathbf{u}(t))\) tends uniformly on K to \(\mathbf{f}(t, \mathbf{w}(t))\) along G, w satisfies equation (6) of IV, p.~165, so is a solution of (1). The uniqueness of the solution follows immediately from inequality (15) of IV, p.~170 where one takes \(\varepsilon_{1} = \varepsilon_{2} = 0\) and \(\mathbf{u}(t_{0}) = \mathbf{v}(t_{0})\) .

We shall say that a function \(\mathbf{f}\) defined on \(\mathrm{I} \times \mathrm{H}\) is locally Lipschitz if, for every point \((t, \mathbf{x})\) of \(\mathrm{I} \times \mathrm{H}\) , there exists a neighbourhood \(\mathrm{V}\) of \(t\) (with respect to \(\mathrm{I}\) ) and a neighbourhood \(\mathrm{S}\) of \(\mathbf{x}\) such that \(\mathbf{f}\) is Lipschitz on \(\mathrm{V} \times \mathrm{S}\) (for a constant \(k\) depending on \(\mathrm{V}\) and \(\mathrm{S}\) ). By the Borel-Lebesgue theorem, for every compact interval \(\mathrm{J} \subset \mathrm{I}\) and every point \(\mathbf{x}_0 \in \mathrm{H}\) there exists an open ball \(\mathrm{S}\) with centre \(\mathbf{x}_0\) , contained in \(\mathrm{H}\) , such that \(\mathbf{f}\) is Lipschitz on \(\mathrm{J} \times \mathrm{S}\) ; thus \(\mathbf{f}\) satisfies the hypotheses of lemma 1 of IV, p.~3. When \(\mathbf{f}\) is locally Lipschitz on \(\mathrm{I} \times \mathrm{H}\) we shall say that the equation \(\mathbf{x}' = \mathbf{f}(t, \mathbf{x})\) is locally Lipschitz on \(\mathrm{I} \times \mathrm{H}\) .

We shall generalize and clarify th. 1 of IV, p.~10 for locally Lipschitz equations; we restrict ourselves to the case where \(t_0\) is the left endpoint of the interval I; one can pass easily to the case where \(t_0\) is an arbitrary point of I (cf.~IV, §IV ??, p.~9, corollary).

THEOREM 2. Let \(\mathbf{I} \subset \mathbf{R}\) be an interval (not reducing to a single point) with left-hand endpoint \(t_0 \in \mathbf{I}\) , let \(\mathbf{H}\) be a nonempty open set in \(\mathbf{E}\) , and \(\mathbf{f}\) a locally Lipschitz function on \(\mathbf{I} \times \mathbf{H}\) .

\(1^{\circ}\) For \(\mathbf{x}_0\in \mathrm{H}\) there exists a largest interval \(\mathbf{J}\subset \mathbf{I}\) , with left-hand endpoint \(t_0\in \mathbf{J}\) , on which there exists an integral \(\mathbf{u}\) of the equation \(\mathbf{x}' = \mathbf{f}(t,\mathbf{x})\) taking values in H and equal to \(\mathbf{x}_0\) at the point \(t_0\) ; this integral is unique.

\(2^{\circ}\) If \(\mathrm{J} \neq \mathrm{I}\) then \(\mathrm{J}\) is a half-open interval \([t_0, \beta[\) of finite length; further, for every compact subset \(\mathbf{K} \subset \mathbf{H}\) the set \(\overline{\mathbf{u}}^{-1}(\mathbf{K})\) is a compact subset of \(\mathbf{R}\) .

\(3^{\circ}\) If \(\mathbf{J}\) is bounded, and if \(\mathbf{f}(t,\mathbf{u}(t))\) is bounded on \(\mathbf{J}\) , then \(\mathbf{u}(t)\) has a left limit \(\mathbf{c}\) at the right-hand endpoint of \(\mathbf{J}\) ; further, if \(\mathbf{J} \neq \mathbf{I}\) then \(\mathbf{c}\) is a boundary point of \(\mathbf{H}\) in \(\mathbf{E}\) .

\(1^{\circ}\) Let \(\mathfrak{M}\) be the set of intervals L (not reducing to a single point) with left-hand endpoint \(t_0 \in L\) which are contained in I and are such that on L there is a solution of (1) (IV, p.~163) with values in H and equal to \(\mathbf{x}_0\) at \(t_0\) ; by th. 1 (IV, p.~171) the set \(\mathfrak{M}\) is not empty. Let L and L' be two intervals belonging to \(\mathfrak{M}\) , and suppose, for example, that L \(\subset\) L'; if u and v are two integrals of (1) defined respectively on L and L', with values in H, and equal to \(\mathbf{x}_0\) at \(t_0\) , we shall see that v is an extension of u. Indeed, let \(t_1\) be the supremum of the set of \(t \in L\) such that \(\mathbf{u}(s) = \mathbf{v}(s)\) for \(t_0 \leqslant s \leqslant t\) ; we shall show that \(t_1\) is the right-hand endpoint of L. If this were not so, we would have \(\mathbf{u}(t_1) = \mathbf{v}(t_1)\) by continuity, and \(\mathbf{x}_1 = \mathbf{u}(t_1)\) would belong to H; since f is locally Lipschitz, th. 1 shows that there can exist only one integral of (1) defined on a neighbourhood of \(t_1\) with values in H and equal to \(\mathbf{x}_1\) at \(t_1\) ; it is therefore a contradiction to suppose that \(t_1\) is not the right-hand endpoint of L. We now see that if J is the union of the intervals L \(\in\) M there exists one and only one integral u of (1), defined on J, with values in H and equal to \(\mathbf{x}_0\) at \(t_0\) .

\(2^{\circ}\) Suppose that \(\mathrm{J} \neq \mathrm{I}\) and let \(\beta\) be the right endpoint of \(\mathrm{J}\) ; if \(\beta\) is the right-hand endpoint of \(\mathrm{I}\) then \(\beta \in \mathrm{I}\) (so \(\beta\) is finite) and \(\mathrm{J} = [t_0, \beta[\) by hypothesis. Suppose then that \(\beta\) is not the right-hand endpoint of \(\mathrm{I}\) ; if \(\beta \in \mathrm{J}\) then \(\mathbf{u}(\beta) = \mathbf{c}\) belongs to \(\mathrm{H}\) ; by th. 1 there exists an integral of (1) with values in \(\mathrm{H}\) , defined on an interval

\[
[ \beta , \beta_ {1} [ \subset I
\]

and equal to \(\mathbf{c}\) at \(\beta\) ; then \(\mathrm{J}\) would not be the largest of the intervals in \(\mathfrak{M}\) , which is absurd; so \(\mathrm{J} = [t_0, \beta[\) .

If K is a compact subset of H then \(\overline{\mathbf{u}}^{-1}(\mathbf{K})\) is closed in J; we shall see that there exists a \(\gamma\in J\) such that \(\overline{\mathbf{u}}^{-1}(\mathbf{K})\) is contained in \([t_{0},\gamma]\) , which will prove that \(\overline{\mathbf{u}}^{-1}(\mathbf{K})\) is compact. If not, there would be a point \(c\in K\) such that \((\beta,\mathbf{c})\) is a cluster point of the set of points \((t,\mathbf{u}(t))\) such that \(t<\beta\) and \(\mathbf{u}(t)\in\mathbf{K}\) . Since \(\beta\in I\) and \(c\in H\) there exists a neighbourhood V of \(\beta\) in I, and an open ball S with centre c and radius r contained in H, such that f is Lipschitz and bounded on V×S; let M be the supremum of \(\|\mathbf{f}(t,\mathbf{x})\|\) over this set. By hypothesis there exists a \(t_{1}\in J\) such that \(\beta-t_{1}<r/2M\) , \(t_{1}\in V\) and \(\|\mathbf{u}(t_{1})-\mathbf{c}\|\leqslant r/2\) ; th. 1 shows that there exists one and only one integral of (1), with values in H, defined on an interval \([t_{1},t_{2}]\) containing \(\beta\) , and equal to \(\mathbf{u}(t_{1})\) at \(t_{1}\) ; since this interval coincides with u on the interval \([t_{1},\beta[\) it follows that \(J=[t_{0},\beta[\) is not the largest interval in M, which is absurd.

\(3^{\circ}\) Suppose that \(\mathbf{J}\) is bounded and that \(\| \mathbf{f}(t,\mathbf{u}(t))\| \leqslant M\) on \(\mathbf{J}\) ; then \(\| \mathbf{u}'(t)\| \leqslant \mathbf{M}\) on the complement of a countable subset of \(\mathbf{J}\) ; then, by the mean value theorem, \(\| \mathbf{u}(s) - \mathbf{u}(t)\| \leqslant \mathbf{M} |s - t|\) for any \(s\) and \(t\) in \(\mathbf{J}\) ; by the Cauchy criterion, \(\mathbf{u}\) has a left limit \(\mathbf{c}\) at the right endpoint \(\beta\) of \(\mathbf{J}\) . If \(\mathbf{J} \neq \mathbf{I}\) then \(\mathbf{c}\) cannot belong to \(\mathbf{H}\) , for on extending \(\mathbf{u}\) by continuity at \(\beta\) , \(\mathbf{u}\) would be an integral of (1) with values in \(\mathbf{H}\) , defined on an interval \([t_0, \beta]\) and equal to \(\mathbf{x}_0\) at \(t_0\) ; then one would have \(\mathbf{J} = [t_0, \beta]\) , contradicting what we have seen in \(2^{\circ}\) .

COROLLARY 1. If H = E and J ≠ I then \(\mathbf{f}(t, \mathbf{u}(t))\) is not bounded on J; if, further, E is finite dimensional, then \(\| \mathbf{u}(t) \|\) has left limit \(+\infty\) at the right-hand endpoint of J.

The first part is an immediate consequence of the third part of th. 2. If E is finite dimensional then every closed ball \(S \subset E\) is compact, so the second part of th. 2 shows that there exists a \(\gamma \in J\) such that \(\mathbf{u}(t) \notin \mathbf{S}\) for \(t > \gamma\) .

If E is finite dimensional it can happen that \(J \neq I\) but that \(\|\mathbf{u}(t)\|\) remains bounded as t tends to the right-hand endpoint of J (IV, p.~200, exerc. 5).

COROLLARY 2. If, on I × H, the function f is Lipschitz with respect to a regulated function k(t), and if the right-hand endpoint β of J belongs to I, then u has a left limit at β; if H = E and if f is Lipschitz with respect to a regulated function k(t) on I × E, then J = I.

Indeed, if \(\beta \in \mathbf{I}\) , there exists a compact neighbourhood \(\mathrm{V}\) of \(\beta\) in \(\mathbf{I}\) , such that \(\mathbf{f}(t, \mathbf{x}_0)\) and \(k(t)\) are bounded on \(\mathrm{V}\) ; then \(\| \mathbf{f}(t, \mathbf{x}) \| \leqslant m \| x \| + h\) ( \(m\) and \(h\) constant)

on \(V \times H\) , whence \(\left\|\mathbf{u}^{\prime}(t)\right\| \leqslant m \left\|\mathbf{u}(t)\right\| + h\) on the complement of a countable subset of \(V \cap J\) , so that \(\left\|\mathbf{u}(t)\right\| \leqslant m \int_{t_{0}}^{t} \left\|\mathbf{u}(s)\right\| ds + q\) (q constant) on \(V \cap J\) ; lemma 2 (IV, p.~170) shows that \(\left\|\mathbf{u}(t)\right\| \leqslant c e^{mt} + d\) (c and d constant) on \(V \cap J\) , and thus \(\mathbf{f}(t, \mathbf{u}(t))\) remains bounded on J, and the corollary results from th. 2 of IV, p.~172.

Examples. 1) For a differential equation of the form \(\mathbf{x}' = \mathbf{g}(t)\) , where \(\mathbf{g}\) is regulated on I, every integral \(\mathbf{u}\) is clearly defined on all of I. One should note that \(\mathbf{u}\) can be bounded on I without \(\mathbf{g}(t)\) being so.

\begin{enumerate}
\def\labelenumi{\arabic{enumi})}
\setcounter{enumi}{1}
\item
  For the scalar equation \(x' = \sqrt{1 - x^2}\) one has \(I = R\) , \(H = ] - 1, +1[\) . If one takes \(t_0 = x_0 = 0\) the corresponding integral is \(\sin t\) on the largest interval containing 0 where the derivative of \(\sin t\) is positive, that is to say, on \(] - \pi / 2, +\pi / 2[\) ; at the endpoints of this interval the integral tends to an endpoint of \(H\) .
\item
  For the scalar equation \(x' = 1 + x^{2}\) one has I = H = R; the integral that vanishes at t = 0 is \(\tan t\) , and the largest interval containing 0 where this function is continuous is \(J = ] - \pi/2, +\pi/2[\) ; and \(|\tan t|\) tends to \(+\infty\) at the endpoints of J (cf.~IV, p.~173, cor. 1).
\item
  For the scalar equation \(x' = \sin tx\) one has I = H = R and the right-hand side is bounded on I × H, so (IV, p.~173, cor. 1) every integral is defined on all of R.
\end{enumerate}

